\documentclass{article}

\usepackage{amsmath}
\usepackage{tabularx}
\usepackage{booktabs}
\usepackage{hyperref}

\title{Performance Report --- \progname}

\author{\authname}

\date{}

\input{../../Comments.text}
\input{../../Common.text}

\begin{document}

\maketitle

\begin{table}[hp]
\caption{Revision History} \label{TblRevisionHistory}
\begin{tabularx}{\textwidth}{llX}
\toprule
\textbf{Date} & \textbf{Developer(s)} & \textbf{Change}\\
\midrule
Apr 20, 2026 & Mirhoseininejad & Initial Performance Report: carries over the
dynamic-benchmark results from VnV Report \S4.1, adds methodology and
improvement log, and records the two performance-driven design decisions
taken during the CAS~741 cycle.\\
\bottomrule
\end{tabularx}
\end{table}

\newpage

\section{Introduction}

\progname{} is a real-time pose-based exercise-feedback application. The
user-facing performance requirement is implicit rather than numerical:
the system must feel responsive while driving an AR overlay rendered on
the live webcam feed. This report makes that requirement concrete by
identifying the performance aspects that matter, the method used to
measure them, the measured results, and the two performance-driven design
changes made during the project. Some of this material is already present
in the VnV Report; it is consolidated here so that a performance-focused
reader does not have to reconstruct the picture from across several
documents.

\section{Performance Aspects}

The aspects of performance that matter for \progname{} are, in priority
order:

\begin{enumerate}
  \item \textbf{Per-frame end-to-end latency.} Time from a webcam frame
    arriving at the browser to that frame's overlay and rep counter
    being rendered. This is the aspect the user notices as ``laggy.''
  \item \textbf{Sustained throughput.} Effective frames per second
    (FPS) sustained under continuous operation. A target of
    $\ge 24$~FPS was adopted as the minimum that still feels smooth to a
    user.
  \item \textbf{Angular-signal stability.} Frame-to-frame jitter of the
    computed joint angle, reported as
    $\overline{|\Delta\theta|}$ (see VnV Report \S4.1 for the precise
    definition). Although this is principally a \emph{quality} signal,
    it also bounds how much smoothing the downstream Kalman filter
    must apply, which in turn affects perceived responsiveness.
  \item \textbf{Test-suite runtime.} The full pytest suite must
    complete quickly to support frequent local iteration.
\end{enumerate}

Two aspects are explicitly out of scope and are recorded here so the
reader knows they were not measured: (i)~resident memory footprint
of the backend process, and (ii)~GPU utilisation, since the production
path runs on CPU only.

\section{Measurement Methodology}

\subsection{Hardware and software under test}

All measurements in this report were taken on the author's macOS
workstation (Apple Silicon; the MediaPipe wheel used is
\texttt{mediapipe-silicon}). The backend runs under Python~3.9.6 with
the package versions pinned in
\texttt{src/backend/requirements.txt}; the frontend runs under Node.js
via Vite. No machine-specific tuning or GPU acceleration was used; the
numbers should therefore be read as a lower-bound baseline for a recent
laptop-class machine.

\subsection{Benchmark drivers}

Three instruments were used:

\begin{itemize}
  \item \textbf{Dynamic benchmark script} (offline). Replays a pre-recorded
    video dataset through each of the three pose backends (MediaPipe 2D,
    MediaPipe 3D, MoveNet 3D) and records both per-frame processing
    latency and the resulting joint-angle time series. The recorded
    dataset consists of short squat and bicep-curl clips captured at
    two camera pitches (eye-level, $45^{\circ}$-down) and five camera
    orientations (front, front-$45^{\circ}$, side, back-$45^{\circ}$,
    back).
  \item \textbf{Live FPS counter} (online). The frontend samples
    \texttt{performance.now()} in its WebSocket render loop and
    reports an exponentially smoothed FPS value. This is the measure
    used to confirm the sustained-throughput target during the POC and
    EXPO demonstrations.
  \item \textbf{pytest runtime.} \texttt{make test} reports wall time.
    The suite consists of the 11 tests described in VnV Report \S6.
\end{itemize}

\subsection{Metrics}

\begin{itemize}
  \item \textbf{Latency (ms):} average per-frame processing time
    reported by the dynamic benchmark script.
  \item \textbf{Stability:}
    $\overline{|\Delta\theta|}
      = \frac{1}{N{-}1}\sum_{i=1}^{N-1}|\theta_i - \theta_{i-1}|$
    on the computed joint-angle signal; lower is more stable.
  \item \textbf{Throughput (FPS):} inverse of end-to-end latency under
    sustained load.
\end{itemize}

\section{Results}

\subsection{Backend comparison}

\begin{table}[h]
\centering
\begin{tabularx}{\textwidth}{lccc}
\toprule
\textbf{Backend} & \textbf{Avg.\ Latency (ms)} &
  \textbf{Knee $\overline{|\Delta\theta|}$ ($^\circ$)} &
  \textbf{Elbow $\overline{|\Delta\theta|}$ ($^\circ$)}\\
\midrule
\texttt{mediapipe\_2d} & ${\sim}39$ & ${\sim}0.7$ & ${\sim}3.8$\\
\texttt{mediapipe\_3d} & ${\sim}41$ & ${\sim}1.2$ & ${\sim}2.3$\\
\texttt{movenet\_3d}   & ${\sim}46$ & ${\sim}1.4$ & ${\sim}5.4$\\
\bottomrule
\end{tabularx}
\caption{Backend comparison on the dynamic benchmark.
$\overline{|\Delta\theta|}$ is a stability measure (lower is more
stable); it is \emph{not} error vs.\ ground truth.}
\label{tab:bench}
\end{table}

MediaPipe 2D is the best choice for this application under every metric
that matters except elbow stability, on which MediaPipe 3D wins by
about $1.5^\circ$. The 2D backend was chosen as the production default.

\subsection{Throughput}

An average per-frame latency of $\sim 39$~ms yields a theoretical
ceiling of $\sim 25.6$~FPS on the production backend. Under live
operation during the POC demonstration (2026-03-05) and the EXPO, the
frontend FPS counter remained at $\ge 24$~FPS for the duration of
each rep set, which clears the target.

\subsection{Test-suite runtime}

\texttt{make test} completes in under 5 seconds locally for the 11
unit tests. This is well within the threshold needed for tight iteration.

\section{Performance Improvements}

Two changes during the project meaningfully improved performance
relative to the CAS~772 monolithic prototype from which \progname{} was
refactored:

\subsection{Backend selection (M4)}

The CAS~772 prototype used only MediaPipe 3D (world landmarks enabled).
Switching to MediaPipe 2D for production reduced per-frame latency from
$\sim 41$~ms to $\sim 39$~ms \emph{and} improved knee-angle stability
from $\sim 1.2^\circ$ to $\sim 0.7^\circ$. The change was small but
cheap: it required only a configuration switch at the M4
(\texttt{m4\_pose\_tracking.py}) boundary, because module M4 hides the
backend behind a common interface. This is a concrete instance of the
modular-decomposition design paying off.

\subsection{Kalman smoothing (M8)}

The raw per-frame angle signal, before smoothing, has a
per-sample standard deviation of $\sigma_{\text{raw}} \approx 5^{\circ}$
on the T\_M8 synthetic noise case. After the Kalman smoother is
applied, the post-transient standard deviation falls below
$\sigma_{\text{raw}}$ in every run. The practical effect on the user is
that the rep-count decision becomes robust to a single noisy frame,
so the state machine does not fire spurious counts.

\section{Bottleneck Analysis}

With the production configuration (MediaPipe 2D + Kalman smoothing on
top of the vectorised angle computation) the dominant cost per frame
is the pose backend itself: approximately $35$--$37$~ms of the
$\sim 39$~ms total budget is spent inside the MediaPipe wheel. The
downstream analytical pipeline (vector math in M5 and the FSM step in
M6) contributes a small, roughly constant overhead of $\sim 2$~ms.
This means further performance work on the analytical pipeline has
diminishing returns; future performance work should target the pose
backend itself (for example, by quantisation, by using the
\texttt{mediapipe} GPU-delegate path on supported hardware, or by
moving to a smaller model such as MoveNet Lightning if a small loss
of stability is acceptable).

\section{Conclusion}

The performance requirements for \progname{} are met by the
production configuration. The sustained throughput of $\ge 24$~FPS and
the knee-angle stability of $\sim 0.7^\circ$ are both adequate for the
intended user experience. The principal residual risk is not a
performance bottleneck --- it is the absence of a physical goniometer
validation of the pose backend against an independent ground truth,
which is documented as open work in the VnV Report (\S4.1).

\end{document}
